% Imporditud: tools/import_eksperimendid.py
% Allikas: Eesti füüsikaolümpiaadi eksperimendi ülesannete
%   kogu (Rael Kalda, 2023), ülesanne Ü47, lahendus L47.
\setAuthor{EFO žürii}
\setRound{lõppvoor}
\setYear{2004}
\setNumber{P E2}
\setDifficulty{2}
\setTopic{vedelike mehaanika}
\setVahendid{anum veega, anum tundmatu vedelikuga, joogitops, joonlaud, marker}
\setPunktid{12}
\setAllikas{Eksperimendi kogu Ü47, lahendus lk 54}
\setLahendusseis{toores}

\prob{Vedeliku tihedus}
Määrake rohelise vedeliku tihedus. Vee tihedus $\rho$ = 1000 kg/m$^3$.
\textit{Märkus:} Sama ülesanne oli ka aastal 2007 põhikooli lõppvoorus.

\hint

\solu
Keha ujub vedelikus, kui kehale mõjuv raskusjõud ja üleslükkejõud on suuruselt võrdsed. Ülelükkejõud võrdub keha poolt välja tõrjutud vedeliku raskusjõuga. Seega, kui keha ujub vees ning tundmatus vedelikus, kehtivad seosed:

mg = $\rho$vgVv, mg = $\rho$gV,

kus indeksiga v on tähistatud olukord, mil keha ujub vees, indeksita aga keha ujumine tundmatus vedelikus. m elimineerimisel saame $\rho$vVv = $\rho$V , millest $\rho$ = $\rho$vVv/V . Üleslükkejõu puhul arvestatakse keha veealuse osa ruumala, mis võrdub keha poolt väljatõrjutud vedeliku ruumalaga. Viimase saab arvutada vedeliku nivoo stabiilselt vedelikus. Mõõdame markerit kasutades veenivoo muutuse topsi vette asetamisel. Mõõdame anuma sisemise läbimõõdu. Kordame sama tundmatu vedeliku korral. Arvutame ruumalad. Arvutame tundmatu vedeliku tiheduse.
\probend
